\section*{Author Section 4: preparations}
\paragraph{Fourier notation.} $\chi_T$ is the Fourier character for the edge-set $T$ in the random graph. The author's Fourier orthonormal basis in Equation(7.1) is $\chi_T=\prod_{e\in T}(x_e-(2p-1))/(2\sqrt{p(1-p)})$ for $G(n,p)$; at $p=1/2$ this is the product of the independent signed edge indicators. This is the standard Fourier convention used throughout the paper.
\paragraph{Remark 3.14 (continued).} In (3.17), the ``most significant'' factor is $(\omega/n)^{|V(T)\cup A\cup B|}\omega^{-c}$, if notice
\[\frac{\binom{n-|V(T)\cup A\cup B|}{u-c}}{\binom{\omega-d+u}u}\omega^un^{-(u-c)}\]
has 0th-order in $\omega,n$. Other factors like $(|V(T)\cup A\cup B|+u-c+8\tau^2)!/(8\tau^2)!$ are qualitatively smaller than $\omega$, within our parameter regime.
\subsection*{4.1 Homogenization for Exact Clique}
With Size Constraints (1.4) satisfied, any moment matrix can be reduced to its $\binom{[n]}{d/2}$-principal minor, which is slightly more convenient. The homogeneity trick is standard in SOS literature. Given a degree-$d$ moment matrix $M^d_{\mathrm{SOS}}(G)$ satisfying (1.4), let $M(G)$ be its principal minor on $\binom{[n]}{d/2}\times\binom{[n]}{d/2}$.
\paragraph{Lemma 4.1.} $M^d_{\mathrm{SOS}}(G)$ is PSD iff $M(G)$ is PSD.
\begin{proof}The $\Rightarrow$ part is trivial. Now suppose $M^d_{\mathrm{SOS}}$ is not PSD, then
\[\exists a\in\mathbb R^{\binom{[n]}{\le d/2}}\quad a^\top M^d_{\mathrm{SOS}}a=-1.\tag{4.1}\]
With boolean constraints (define $\widetilde{\mathbb E}(x_i^2p):=\widetilde{\mathbb E}(x_ip)$ for all $i$ and polynomial $p$ of degree $\le d-2$), this is equivalent to $\widetilde{\mathbb E}(g^2)=-1$ (4.2), where $g=a^\top x=\sum_{|S|\le d/2}a_Sx_S$ is multi-linear. Now substitute every $x_S$ ($|S|<d$) in $g$ by the corresponding linear combination of $\{x_{S'}:|S'|=d\}$ from (3.13). This does not affect (4.2) since $\widetilde{\mathbb E}$ satisfies the equality constraints. We get $\widetilde{\mathbb E}(g_1^2)=-1$ (4.3) for some multi-linear, degree-$d/2$ homogeneous $g_1$. Translate (4.3) back (assume $g_1=b^\top x$, $x=(x_S)_{|S|=d/2}$) to $b^\top Mb=-1$, so $M$ is not PSD.\end{proof}
\subsection*{4.2 Concentration bound on polynomials}
\paragraph{Lemma 4.2.} Suppose $a<\log n$, and $p=\sum_{T:|V(T)|=a}c(T)\chi_T$, $c_T\in\mathbb R$, is a polynomial, and $C>0$ satisfies $|c(T)|\le C$ for all $T$. Then w.p. $1-n^{-10\log n}$ over $G$,
\[|p(G)|<C\,n^{a/2}2^{a^2}n^{4\log\log n}.\tag{4.4}\]
\begin{proof}Power-estimation. For all $k\in\mathbb N$ (think of $a<k=o(n/a)$),
\[p^{2k}=\sum_{T_1,\ldots,T_{2k}:|V(T_i)|=a}c(T_1)\cdots c(T_{2k})\chi_{T_1}\cdots\chi_{T_{2k}}.\tag{4.5}\]
Take expectation. Each $\mathbb E[\chi_{T_1}\cdots\chi_{T_{2k}}(G)]\ne0$ (equals 1) iff every edge appears even times in $T_1,\ldots,T_{2k}$, which implies $|V(T_1\cup\cdots\cup T_{2k})|\le\frac12\cdot2ka=ka$. There are at most $\binom n{ka}<n^{ka}$ choices of such $V(T_1\cup\cdots\cup T_{2k})$. For each choice, at most $\binom{ka}a2^{\binom a2}<(ka)^a2^{a^2/2}$ ways to choose each $T_i$. Therefore,
\[\mathbb E[p^{2k}]\le C^{2k}n^{ka}((ka)^a2^{a^2/2})^{2k}=:N^{2k},\quad N=Cn^{a/2}(ka)^a2^{a^2/2}.\]
By Markov, $\Pr[p^{2k}>(2N)^{2k}]<2^{-2k}$. Take $k:=10\log^2n$; w.p. $>1-n^{-10\log n}$, $|p(G)|<2N<Cn^{a/2}2^{a^2}n^{4\log\log n}$ for all large enough $n$.\end{proof}
\subsection*{4.3 Norm concentration of pseudo-random matrices}
Pseudo-random matrices refer to graphical matrices [24]: intuitively, such a matrix collects Fourier characters of all embeddings of a fixed ``shape''. Definitions4.3/4.5 are implicit in [24,25,16] and termed explicitly in [5].
\paragraph{Definition 4.3.} A ribbon $R$ is an ordered triple $(A,B;T)$ with $A,B$ vertex-sets and $T$ an edge-set. $A,B$ are its left/right vertex-sets. Its size is $|V(R)|=|V(T)\cup A\cup B|$. By definition a ribbon as a graph has no isolated vertex outside $A\cup B$.
\paragraph{Definition 4.4.} $R=(A,B;T)$ is left-generated if every vertex in $V(R)$ is either in $B$ or can be reached by paths from $A$ without touching $B$.\footnote{The author always includes degenerate paths (one-point paths).} Right-generated is symmetric.
\paragraph{Definition 4.5.} A ribbon $(A,B;T)$ is a shape if $A\cup B$ is totally ordered. Denote a shape by $U=(A,B;T)$, with $V(U)=A\cup B\cup V(T)$, size $|V(U)|$. When fixing underlying vertex-set $[n]$, a ribbon within $[n]$ can always be regarded as shapes with the induced vertex ordering. The author may interchangeably use $R$ to denote shapes.
\paragraph{Definition 4.6.} A real-valued function $f$ on ribbons within $[n]$ is symmetric with respect to shapes if $f(R)=f(R')$ whenever $R,R'$ have the same shape.
\paragraph{Definition 4.7 ([24]).} Fix $n$ and shape $U=(A,B;T)$. A map $\phi:V(U)\to[n]$ is proper if injective and respects the order on $A\cup B$. Its graphical matrix is the $2^{[n]}\times2^{[n]}$ matrix
\[\forall I,J\subseteq[n],\quad M_U(I,J)=\sum_{T':\exists\text{ proper }\phi,\ \phi(A)=I,\phi(B)=J,\phi(T)=T'}\chi_{T'}\]
(0 if no such $\phi$ exists). Here $\phi$ on $T$ is the natural induced edge map.
\paragraph{Theorem 4.8 (Norm bounds on $M_U$ [24,5]).} For any shape $U=(A,B;T)$ of size $t<\log n$, w.p. $>1-n^{-10\log n}$ over $G$,
\[\|M_U(G)\|\le n^{(t-p)/2}2^{O(t)}(\log n)^{O(t+p-2r)},\tag{4.6}\]
where $r=|A\cap B|$, $p$ is the maximum number of vertex-disjoint paths between $(A,B)$ in $U$. This is tight up to $\mathrm{polylog}(n)$ factors for all $M_U$ with these parameters ([24],Theorem38). Moreover, if $s=(|A|+|B|)/2$, then
\[\|M_U(G)\|\le n^{(t-p)/2}2^{O(t)}(\log n)^{O(t-s)}.\tag{4.7}\]
The author cites the standard trace-power norm theorem without reproducing its earlier proof. The ``moreover'' part follows since $t\ge|A\cup B|=2s-r$, $p\le s$, so $t+p-2r\le t+s-2(2s-t)=3(t-s)$.
\paragraph{Remark 4.9.} Shapes $U=(A,A;T)$ give diagonal $M_U$, so estimating its norm is estimating the absolute values of its polynomial diagonals; Theorem4.8 generalizes Lemma4.2.
\subsection*{4.4 General notions on graphs}
\paragraph{Definition 4.10 (Vertex-separator).} Given graph $H$ and $A,B\subseteq V(H)$, call $S\subseteq V(H)$ an $(A,B)$-vertex-separator if every path from $A$ to $B$ in $H$ must pass through $S$.\footnote{Degenerate paths are included, as before.} Let $s_{A,B}(H):=\min\{|S|:S\text{ is an }(A,B)\text{-vertex-separator}\}$. A separator achieving this is a min-separator; $\mathrm{mSep}_{A,B}(H)$ is the set of all min-separators. For ribbon $R=(A,B;T)$ use the graph on $V(T)\cup A\cup B$ with edges $T$ and write $s_{A,B}(T)$, $\mathrm{mSep}_{A,B}(T)$ or $\mathrm{mSep}(R)$.
\paragraph{Menger's theorem.} For finite graph $H$, $s_{A,B}(H)$ equals the maximum number of vertex-disjoint paths from $A$ to $B$ in $H$.
\paragraph{Definition 4.11.} The reduced size of ribbon $R=(A,B;T)$ is\[e_{A,B}(T):=|V(T)\cup A\cup B|-s_{A,B}(T).\tag{4.8}\]
It is double the exponent of $n$ in Theorem4.8's bound, hence controls the norm of the graphical matrix.
